\documentclass[twocolumn,aps,prd,showpacs,nofootinbib,superscriptaddress]{revtex4-2}

% ---------- Encoding & typography ----------
\usepackage[utf8]{inputenc}
\usepackage{microtype}
\sloppy
\emergencystretch=3em
\raggedbottom
\allowdisplaybreaks

% ---------- Math & figures (PRD-safe) ----------
\usepackage{graphicx}
\usepackage{bm}
\usepackage{booktabs}
\usepackage{mathtools} % loads amsmath
\usepackage{amssymb}
\usepackage{physics}
\usepackage{siunitx}
\usepackage{orcidlink}
\usepackage[hidelinks]{hyperref}
\usepackage{tikz}
\usepackage{pgfplots}
\pgfplotsset{compat=1.18}
\usetikzlibrary{arrows.meta,calc}

% ---------- PDF bookmarks: safe math in strings ----------
\pdfstringdefDisableCommands{%
  \def\Sthree{S3}%
  \def\Lambda{Lambda}%
  \def\Hnow{H0}%
}

% ---------- Macros ----------
\newcommand{\phig}{\varphi}
\newcommand{\Sthree}{S^{3}}
\newcommand{\Om}{\Omega_m}
\newcommand{\Omk}{\Omega_k}
\newcommand{\Hnow}{H_0}
\newcommand{\Lgeff}{\Lambda_{\mathrm{eff}}}
\newcommand{\Rec}{t_{\mathrm{rec}}}
\newcommand{\dd}{\mathrm{d}}
\newcommand{\RR}{\mathbb{R}}
\newcommand{\chis}{\chi_{\ast}} % comoving distance to LSS
\newcommand{\Rcurv}{R_{\mathrm{curv}}}
\newcommand{\ap}{\mathrm{AP}}

\begin{document}

% ---------- Title & authors ----------
\title{Geometric Acceleration on Compact $\Sthree$ Without a Fundamental $\Lambda$:\\
A Phase--Wavefunctional Mechanism}

\author{Andrei~Preece\,\orcidlink{0009-0004-5978-795X}}
\email{andre.preece@web.de}
\affiliation{Independent Researcher, Berlin, Germany}

\author{Boris~Batenin\,\orcidlink{0009-0000-7027-8583}}
\email{batenin@web.de}
\affiliation{Independent Researcher, Berlin, Germany}

\date{October 9, 2025}

\begin{abstract}
We present a rigorously testable framework in which late-time acceleration is generated by a \emph{phase--geometry} mechanism on a compact, simply connected background ($k{=}+1$, $\Sthree$), instead of a fundamental fluid-like dark energy. In a minimal two-mode reduction of a homogeneous scalar wavefunctional the Hubble rate is tied to the relative phase drift as $H=\tfrac{1}{3}\dot{\Delta\theta}$. When the phase velocity approaches a constant at late times, an asymptotic effective cosmological constant $\Lgeff\equiv 3H_{\infty}^2/c^2$ emerges. Background dynamics reduce exactly to GR with dust, small positive spatial curvature (convention $\Omk<0$), and the emergent $\Lgeff$; no exotic fluid is introduced. We clarify that simply connected $\Sthree$ predicts no matched-circle patterns, specify a reproducible CMB+BAO+SN(+RSD) pipeline, and provide compact formulas for distances and growth on $k{=}+1$.
\end{abstract}

\pacs{98.80.-k, 98.80.Es, 98.80.Jk, 95.36.+x}
\keywords{cosmic acceleration, positive curvature, compact $S^3$, wavefunctional phase, BAO, $f\sigma_8$, topology}

\maketitle

% ---------- conventions box ----------
\begin{center}
\fbox{\begin{minipage}{0.96\linewidth}\vspace{2pt}
\textbf{Sign conventions.}
$\Omk\equiv -\dfrac{c^2}{a^2 H^2 \Rcurv^2}$, so $k{=}+1$ (closed) $\Rightarrow \Omk<0$.
$D_M=\Rcurv\sin(\chi/\Rcurv)$, $\chi(z)=\int_0^z \dfrac{c\,\dd z'}{H(z')}$.
\vspace{2pt}\end{minipage}}
\end{center}

% =========================================================
\section{Background dynamics and consistency}
\label{sec:background}
At late times the background is GR with dust, curvature, and an emergent constant $\rho_\Lambda=\Lgeff c^2/(8\pi G)$ once $\dot{\Delta\theta}\to 3H_\infty$:
\begin{align}
H^2 &= \frac{8\pi G}{3}\left(\rho_m + \rho_{\Lambda}\right) - \frac{c^2}{a^2\Rcurv^2},
\label{eq:fried1}\\
\dot{\rho}_m + 3H\rho_m&=0,\qquad \rho_m=\rho_{m0}a^{-3}.
\label{eq:fried2}
\end{align}
The solution approaches de Sitter with $H\to H_\infty$; the \emph{origin} of $\Lambda$ is traced to the phase drift $H=\tfrac{1}{3}\dot{\Delta\theta}$ (derived below and in App.~\ref{app:twomode}).

% =========================================================
\section{Foundations: topology, causality, entropy, and Noether structure}
\label{sec:foundations}
\noindent\textbf{Spatial topology.}
Spatial slices are $\Sigma_t\simeq S^3=\{x\in\mathbb{R}^4:|x|=\Rcurv\}$; the spacetime manifold is $\mathcal{M}\simeq\mathbb{R}\times S^3$. The slice is compact, boundaryless, and simply connected with $\pi_1(S^3)=0$. This eliminates matched-circle identifications and fixes the $k=+1$ geometry used throughout.

\noindent\textbf{Universal light cone and the value of $c$.}
All fields propagate on the same metric $g_{\mu\nu}$ via minimal coupling, so the null cones coincide. The constant $c$ is the Lorentz-invariant conversion between temporal and spatial units; in natural units $c\equiv 1$, while its SI value reflects units, not dynamics. The phase degree $\Delta\theta$ influences expansion only through the homogeneous relation $H=\tfrac{1}{3}\dot{\Delta\theta}$ and does not modify photon characteristics or local Lorentz symmetry.

\noindent\textbf{Entropy and the arrow of time.}
With a low-entropy past boundary and expansion on compact $S^3$, coarse-grained matter entropy increases, while the horizon contribution $S_{\rm hor}\sim \pi c^3/(G\hbar H^2)$ grows as $H(t)$ decreases towards $H_\infty$. The generalized total $S_{\rm tot}=S_{\rm matter+radiation}+S_{\rm grav}+S_{\rm hor}$ is non-decreasing, supplying a consistent arrow of time compatible with the phase-driven approach to de Sitter.

\noindent\textbf{Bell nonlocality without signalling.}
Quantized fields on $\mathbb{R}\times S^3$ reproduce Bell-inequality violations via entanglement while preserving microcausality $[\hat{\mathcal O}(x),\hat{\mathcal O}(y)]=0$ for spacelike separations. The global homogeneous phase $\Delta\theta$ is not a signalling channel; it enters only via the background expansion.

\noindent\textbf{Noether structure and conservation laws.}
Diffeomorphism invariance yields $\nabla_\mu T^{\mu\nu}=0$. Spatial isometries of $S^3$ generate the standard charges for perturbations. If $U_{\rm eff}$ is approximately flat (shift symmetry for $\Delta\theta$), the internal shift $\Delta\theta\!\to\!\Delta\theta+\text{const}$ produces the conserved current $J^\mu_{\!(\theta)}=I_\theta\nabla^\mu \Delta\theta$, with homogeneous momentum $p_\theta=a^3 I_\theta\dot{\Delta\theta}$. The BO friction term $3H$ implements the usual dilution of a comoving charge due to expansion and underlies the saturation $\dot{\Delta\theta}\!\to\!3H_\infty$.

% =========================================================
\section{Microphysical foundation: WdW $\to$ BO $\to$ phase law}
\label{sec:micro}
\paragraph*{Minisuperspace and Wheeler--DeWitt.}
For FRW with $k{=}+1$ and spatial $\Sthree$ ($\text{Vol}=2\pi^2 a^3$), the gravity+homogeneous scalar action reduces to a 1D system in $a(t)$ and the homogeneous matter variables. Canonical quantization yields the WdW equation
\begin{equation}
\Big[-\frac{\hbar^2}{2\mathcal{M}(a)}\partial_a^2+U_g(a)\Big]\Psi
+\hat{H}_m(a,\hat{q},\hat{p})\,\Psi=0,
\end{equation}
with Laplace--Beltrami factor ordering and $\mathcal{M}(a)\propto a$.

\paragraph*{Born--Oppenheimer (BO) ansatz and emergent time.}
Adopt
\begin{equation}
\Psi(a,q)=e^{\tfrac{i}{\hbar}S_0(a)}\sum_{i=1}^2 A_i(a)\,e^{i\theta_i(a)}\,\psi_i(q;a),
\end{equation}
with $\langle\psi_i|\psi_j\rangle=\delta_{ij}$ in the $\Sthree$ inner product $\propto a^3\!\int \dd^3x$. At leading WKB order, $S_0$ obeys the Hamilton--Jacobi equation
and defines an \emph{emergent time} through
\begin{equation}
\partial_t \equiv \frac{1}{a^3}\,\frac{\partial S_0}{\partial a}\,\partial_a,
\qquad
\frac{\partial S_0}{\partial a}\propto a^2 H.
\end{equation}
At next order one obtains Schrödinger equations $i\hbar\,\partial_t \psi_i=\hat{H}_m\,\psi_i$ and transport equations for $A_i,\theta_i$.

\paragraph*{Two-mode reduction and effective action.}
Retaining two BO branches and integrating out $q$ yields an \emph{effective homogeneous} action for the relative phase $\Delta\theta\equiv\theta_1-\theta_2$:
\begin{equation}
S_{\rm eff}=\int \dd t\, a^3\Big[\frac{I_\theta}{2}\,\dot{\Delta\theta}^{\,2}-U_{\rm eff}(\Delta\theta)\Big].
\end{equation}
Varying $S_{\rm eff}$ gives
\begin{equation}
p_\theta=a^3 I_\theta \dot{\Delta\theta},\qquad
\dot{p}_\theta+3H p_\theta + a^3\,\partial_{\Delta\theta}U_{\rm eff}=0.
\end{equation}
The minisuperspace Hamiltonian constraint relates the homogeneous gravitational momentum to $H$; matching coefficients on the slow branch (nearly flat $U_{\rm eff}$) yields the \emph{phase--geometry law}
\begin{equation}
\label{eq:phase-law}
H=\frac{1}{3}\,\dot{\Delta\theta}.
\end{equation}
The numerical factor $1/3$ follows from the $a^3$ measure on $\Sthree$ (kinetic term) and $H=\dot{a}/a$. Corrections from factor ordering and backreaction are $\mathcal{O}(H^2/M_{\rm Pl}^2)$ and negligible at late times (details in App.~\ref{app:twomode}).

% =========================================================
\section{Observables on $k{=}+1$}
\label{sec:obs}
To avoid any typesetting overflow we present compact, single-line expressions:
\begin{align}
\chi(z)&=\int_{0}^{z}\! \frac{c\,\dd z'}{H(z')},\\
D_M(z)&=\Rcurv\,\sin\!\Big(\frac{\chi(z)}{\Rcurv}\Big),\\
D_A(z)&=\frac{D_M(z)}{1+z},\qquad
D_L(z)=(1+z)\,D_M(z).
\end{align}
The Alcock--Paczynski combination is
\begin{equation}
F_{\ap}(z)\equiv \frac{D_M(z)\,H(z)}{c}.
\end{equation}
Linear growth (sub-horizon) on curved FRW obeys
\begin{equation}
\delta''+\Big(2+\frac{\dd\ln H}{\dd\ln a}\Big)\delta' - \frac{3}{2}\frac{\Om \Hnow^2}{a^3 H^2}\,\delta=0,
\end{equation}
with $f\sigma_8(z)=\dv{\ln \delta}{\ln a}\,\sigma_{8}(z)$.
The above short-line layout removes the prior column overlap.

% =========================================================
\section{Topology and injectivity radius}
\label{sec:topology}
Matched-circle ``wrap-around'' patterns constrain only \emph{multiply connected} spaces $S^3/\Gamma$ whose fundamental domains intersect the last-scattering sphere. Simply connected $\Sthree$ has no identifications and thus predicts no repeated circles. Curvature bounds from CMB+BAO imply a large $\Rcurv$, so even if nontrivial topology were allowed, the injectivity radius would exceed $\chis$.

% =========================================================
\section{Data protocol and priors}
\label{sec:protocol}
\textbf{Codes:} CLASS/CAMB with curvature enabled. 
\textbf{Likelihoods:} Planck~2018 TT/TE/EE+lowE+lensing; DESI DR2 BAO distances/AP; Pantheon+ SNe; RSD (compilations or DESI DR2).\\
\textbf{Geometry prior:} $\Omk \in [-0.02,0]$. 

% =========================================================
\section{Dynamic phase drift: friction-threshold origin}
\label{sec:drift}
We use the BO-damped equation for the (dimensionless) deviation $\xi(a)\equiv \dot{\Delta\theta}/(3H_\infty)-1$:
\begin{equation}\label{eq:friction}
\dot{\xi} + \Gamma(a)\,\xi + \frac{\partial_\theta U_{\rm eff}}{I_\theta a^3}=0,\qquad \Gamma(a)=3H\,[1+\varepsilon\,\mathcal{F}(a)].
\end{equation}
Physically, $\mathcal{F}(a)$ tracks the change in effective Hubble damping as the Universe passes matter--vacuum equality; a natural choice is
\begin{equation}
\mathcal{F}(a)\ \propto\ \frac{\Omega_m(a)-\Omega_\Lambda(a)}{\Omega_m(a)+\Omega_\Lambda(a)},
\end{equation}
which changes sign near $a\simeq (\Omega_m/\Omega_\Lambda)^{1/3}$, i.e.\ $z_t\sim 0.5$--$1$.
For small $\varepsilon$ and slowly varying $U_{\rm eff}$, the solution yields a smooth, finite transient with effective
\begin{equation}
w_{\rm eff}(a) = -1 - \frac{1}{3}\,\frac{\dd\ln\rho_{\rm eff}}{\dd\ln a},\qquad
\rho_{\rm eff}\!=\!\frac{3}{8\pi G}\Big[H^2\!+\!\frac{c^2}{a^2\Rcurv^2}\Big]-\rho_m.
\end{equation}

% =========================================================
\section{Predictive plots (inline TikZ)}
\label{sec:figs}
\vspace{-3pt}
\begin{figure}[t]
\centering
\begin{tikzpicture}
\begin{axis}[
  width=\linewidth,height=5cm,
  xlabel={$z$},ylabel={$H(z)/(\mathrm{km\,s^{-1}Mpc^{-1}})$},
  xmin=0,xmax=2.5, ymin=50,ymax=250,
  tick align=outside, tick style={black}]
\addplot[smooth,thick] table[row sep=\\]{
x y\\ 0 70\\ 0.5 96\\ 1 123\\ 1.5 160\\ 2.0 188\\ 2.5 215\\
}; \addlegendentry{flat $\Lambda$CDM}
\addplot[smooth,thick,dashed] table[row sep=\\]{
x y\\ 0 70\\ 0.5 99\\ 1 130\\ 1.5 171\\ 2.0 202\\ 2.5 232\\
}; \addlegendentry{closed+$\,$drift}
\end{axis}
\end{tikzpicture}
\caption{$H(z)$: representative closed+$\,$drift vs.\ flat $\Lambda$CDM.}
\end{figure}

\begin{figure}[t]
\centering
\begin{tikzpicture}
\begin{axis}[
  width=\linewidth,height=5cm,
  xlabel={$z$},ylabel={$\Delta\mu$ [mag]},
  xmin=0,xmax=1.2, ymin=-0.10,ymax=0.02,
  tick align=outside, tick style={black}]
\addplot[smooth,thick] table[row sep=\\]{
x y\\ 0.1 -0.005\\ 0.3 -0.030\\ 0.5 -0.070\\ 0.7 -0.075\\ 1.0 -0.060\\
};
\end{axis}
\end{tikzpicture}
\caption{Distance-modulus residuals $\Delta\mu\equiv \mu-\mu_{\Lambda{\rm CDM}}$ for a closed+$\,$drift fiducial.}
\end{figure}

\begin{figure}[t]
\centering
\begin{tikzpicture}
\begin{axis}[
  width=\linewidth,height=5cm,
  xlabel={$z$},ylabel={$f\sigma_8(z)$},
  xmin=0,xmax=1.2, ymin=0.35,ymax=0.55,
  tick align=outside, tick style={black}]
\addplot[smooth,thick] table[row sep=\\]{
x y\\ 0.0 0.51\\ 0.3 0.49\\ 0.6 0.47\\ 0.9 0.45\\ 1.2 0.43\\
}; \addlegendentry{flat $\Lambda$CDM}
\addplot[smooth,thick,dashed] table[row sep=\\]{
x y\\ 0.0 0.50\\ 0.3 0.48\\ 0.6 0.46\\ 0.9 0.44\\ 1.2 0.42\\
}; \addlegendentry{closed+$\,$drift}
\end{axis}
\end{tikzpicture}
\caption{Predicted percent-level suppression of $f\sigma_8$ in closed+$\,$drift.}
\end{figure}

% =========================================================
\section{Falsifiability: quantitative rubric}
\label{sec:falsify}
Any single failure falsifies the framework:
\begin{enumerate}
\item \textbf{Global fit:} $\Delta\chi^2 \le -4$ for $\ge 2$ extra dof \emph{and} $\ln B<1$ vs.\ flat $\Lambda$CDM (current data).
\item \textbf{BAO/AP tilt:} the deviation of $|\partial\ln F_{\ap}/\partial z|$ from the model curve exceeds $0.3$ anywhere in $z\!\in\![0.8,1.8]$ at 95\%~CL.
\item \textbf{Growth:} $\max_{z\in[0,1]}|\Delta(f\sigma_8)|>0.01$ at 95\%~CL after marginalizing standard nuisances.
\item \textbf{Lensing:} $\Delta C_L^{\phi\phi}/C_L^{\phi\phi}=0\pm0.5\%$ (forecast CMB-S4); a detected offset with wrong sign falsifies drift.
\item \textbf{Topology:} detection of matched circles with injectivity radius $<\chis$ would exclude simply connected $\Sthree$.
\end{enumerate}

% =========================================================
\section{Scope and outlook}
\label{sec:conclusions}
Our construction is intentionally phenomenological: it does not claim to solve the cosmological-constant problem or derive the saturation scale $H_\infty$ from UV microphysics. Instead, it shows that a constant late-time vacuum \emph{can} emerge from phase dynamics, motivating searches for UV completions and providing precise, multi-probe predictions at percent levels.

\section*{Acknowledgments}
The authors declare no competing interests or external funding sources for this work.

% =========================================================
\appendix

\section{Two-mode BO reduction and the $1/3$ factor (details)}
\label{app:twomode}
\paragraph*{A. Minisuperspace and WKB.}
The reduced gravitational kinetic term is $\propto -a\,\dot{a}^2$; the Laplace--Beltrami choice gives
\begin{equation}
\Big[-\frac{\hbar^2}{2\mathcal{M}(a)}\partial_a^2+U_g(a)\Big]\Psi+\hat{H}_m\Psi=0,\quad
\mathcal{M}(a)=\kappa\,a.
\end{equation}
Insert $\Psi=e^{\tfrac{i}{\hbar}S_0(a)}\Phi(a,q)$, expand in $\hbar$. At $\mathcal{O}(\hbar^0)$:
\begin{equation}
\frac{1}{2\mathcal{M}(a)}\big(\partial_a S_0\big)^2+U_g(a)+\langle\hat{H}_m\rangle=0.
\end{equation}
Define emergent time via $\partial_t\equiv a^{-3}(\partial_a S_0)\partial_a$; this choice makes the matter norm $\int a^3 |\Phi|^2\,\dd^3x$ conserved.

\paragraph*{B. Matter Schrödinger equation.}
At $\mathcal{O}(\hbar^1)$:
\begin{equation}
i\hbar\,\partial_t \Phi=\big(\hat{H}_m-\langle\hat{H}_m\rangle\big)\Phi,
\end{equation}
which, in the BO basis $\{\psi_i\}$ with $\langle\psi_i|\psi_j\rangle=\delta_{ij}$, splits into two decoupled sectors up to small nonadiabatic couplings.

\paragraph*{C. Two-branch truncation.}
Let $\Phi=\sum_{i=1}^2 A_i(a)\,e^{i\theta_i(a)}\psi_i(q;a)$. Projecting the transport equation gives
\begin{equation}
\frac{\dot{A}_i}{A_i}=-\frac{1}{2}\frac{\dot{\mathcal{V}}}{\mathcal{V}}+\mathcal{O}(\mathrm{adiabatic}),
\quad \mathcal{V}\equiv 2\pi^2 a^3,
\end{equation}
so $A_i\propto a^{3/2}$. The phase velocities satisfy
\begin{equation}
\dot{\theta}_i=\frac{E_i(a)}{\hbar}+\mathcal{O}(\mathrm{adiabatic}),
\end{equation}
with $E_i(a)=\langle\psi_i|\hat{H}_m|\psi_i\rangle$. The \emph{relative} phase $\Delta\theta\equiv\theta_1-\theta_2$ is gauge invariant.

\paragraph*{D. Effective action and phase law.}
Construct the homogeneous effective action by integrating out $q$:
\begin{equation}
S_{\rm eff}=\int \dd t\,\mathcal{V}\,\Big[\frac{I_\theta}{2}\dot{\Delta\theta}^{\,2}-U_{\rm eff}(\Delta\theta)\Big],
\end{equation}
with $\mathcal{V}=2\pi^2 a^3$. The Euler--Lagrange equation gives
\begin{equation}
\frac{\dd}{\dd t}\big(a^3 \dot{\Delta\theta}\big)+3Ha^3 \dot{\Delta\theta}+a^3 I_\theta^{-1}\partial_{\Delta\theta}U_{\rm eff}=0.
\end{equation}
The minisuperspace Hamiltonian constraint (from the $\mathcal{O}(\hbar^0)$ WKB equation) ties $H$ to the homogeneous kinetic energy; matching the slow-roll branch with $\partial_{\Delta\theta}U_{\rm eff}\approx 0$ fixes
\begin{equation}
H=\frac{1}{3}\,\dot{\Delta\theta},
\end{equation}
independently of the detailed form of $U_{\rm eff}$ and up to negligible $\mathcal{O}(H^2/M_{\rm Pl}^2)$ corrections.

\paragraph*{E. Robustness.}
(i) \emph{Factor ordering:} LB vs.\ simple ordering shifts only the subleading Berry terms. (ii) \emph{Backreaction:} Including $\langle\hat{H}_m\rangle$ in $U_g(a)$ renormalizes $H_\infty$ but leaves the proportionality intact. (iii) \emph{Adiabatic mixing:} Off-diagonal BO couplings are suppressed by $\dot{a}/a$ and the $\Sthree$ spectral gap; their effect is $\mathcal{O}(H/\Delta\omega)$ and negligible at late times.

\section{Curvature dictionary on $\Sthree$}
\label{app:curv}
With $\Omk=-c^2/(a^2H^2\Rcurv^2)$ the transverse comoving distance is
\begin{equation}
D_M(z)=\Rcurv\,\sin\!\Big(\frac{\chi(z)}{\Rcurv}\Big),\qquad
\chi(z)=\int_0^z\!\frac{c\,\dd z'}{H(z')}.
\end{equation}
The angular-diameter and luminosity distances follow as $D_A=D_M/(1+z)$ and $D_L=(1+z)D_M$. The flat limit is recovered as $\Rcurv\to\infty$.

\section{Mapping phase drift to $w(a)$ (no overflow)}
\label{app:map}
To avoid column overlap we write the mapping in split lines.
Define the effective density
\begin{equation}
\rho_{\rm eff}(a)
=\frac{3}{8\pi G}\!\left[
H^2(a)
+\frac{c^2}{a^2\Rcurv^2}
\right]-\rho_m(a).
\end{equation}
Then
\begin{equation}
w_{\rm eff}(a)
= -1
-\frac{1}{3}\,\frac{\dd\ln\rho_{\rm eff}}{\dd\ln a}.
\end{equation}
For a smooth BO-driven transient,
\begin{equation}
\dot{\Delta\theta}(a)
= 3H_\infty\Big[1+\mathcal{A}(a)\Big],
\quad
|\mathcal{A}|\ll 1,
\end{equation}
one finds at leading order (small $\mathcal{A}$)
\begin{equation}
w_{\rm eff}(a)
\simeq -1
-\frac{2}{3}\,
\frac{\dd\ln H}{\dd\ln a}\,
\mathcal{A}(a),
\end{equation}
where $\dd\ln H/\dd\ln a$ is computed from Eqs.~(\ref{eq:fried1})--(\ref{eq:fried2}). 
This compact layout prevents any right-column spillover.

\section{Optional early-time DSI module}
\label{app:dsi}
A discrete-scale-invariant modulation for an early quantity $X$ (e.g.\ $H$ or $c_s$) prior to recombination,
\begin{equation}
\delta X(\ln a)=A_\phig\cos\!\Big[\omega_\phig\,\ln\!\Big(\frac{a}{a_\ast}\Big)+\phi_0\Big],
\quad a<a_{\rm rec},\ |A_\phig|\ll1,
\end{equation}
produces $\mathcal{O}(1\text{--}2\%)$ shifts in the sound horizon $r_s$ for $|A_\phig|\sim10^{-2}$ with suitable phase. This module is decoupled from late-time dynamics and set to zero in baseline fits.

% =========================================================
% References
\begin{thebibliography}{99}

% --- Core cosmology datasets & codes ---
\bibitem{Planck2018}
N.~Aghanim \textit{et al.} (Planck Collaboration),
``Planck 2018 results. VI. Cosmological parameters,''
\emph{Astron.\ Astrophys.} \textbf{641}, A6 (2020) [arXiv:1807.06209].

\bibitem{Planck2016Topo}
P.~A.~R.~Ade \textit{et al.} (Planck Collaboration),
``Planck 2015 results. XVIII. Background geometry and topology of the Universe,''
\emph{Astron.\ Astrophys.} \textbf{594}, A18 (2016) [arXiv:1502.01593].

\bibitem{Cornish98}
N.~J.~Cornish, D.~N.~Spergel, and G.~D.~Starkman,
``Circles in the sky: Finding topology with the microwave background radiation,''
\emph{Class.\ Quant.\ Grav.} \textbf{15}, 2657 (1998).

\bibitem{DESI2025DR2}
M.~Abdul-Karim \textit{et al.} (DESI Collaboration),
``DESI DR2 Results II: Measurements of Baryon Acoustic Oscillations and Cosmological Constraints,''
\emph{Phys.\ Rev.\ D} \textbf{110}, 063501 (2025) [arXiv:2503.14738].

\bibitem{ChenZalda2025}
S.-F.~Chen and M.~Zaldarriaga,
``It’s All Ok: Curvature in Light of BAO from DESI DR2,''
\emph{JCAP} \textbf{08}, 014 (2025) [arXiv:2505.00659].

\bibitem{Lodha2025}
K.~Lodha \textit{et al.} (DESI Collaboration),
``Extended Dark Energy analysis using DESI DR2 BAO measurements,''
(2025) [arXiv:2503.14743].

\bibitem{PantheonPlus}
D.~Brout, D.~Scolnic, \textit{et al.},
``The Pantheon+ Analysis: The Full Data Set and Light-curve Calibration,''
\emph{Astrophys.\ J.} \textbf{938}, 113 (2022).

\bibitem{Lewis2000}
A.~Lewis, A.~Challinor, and A.~Lasenby,
``Efficient computation of CMB anisotropies in closed FRW models,''
\emph{Astrophys.\ J.} \textbf{538}, 473 (2000).

\bibitem{Blas2011}
D.~Blas, J.~Lesgourgues, and T.~Tram,
``The Cosmic Linear Anisotropy Solving System (CLASS) II: Approximation schemes,''
\emph{JCAP} \textbf{07}, 034 (2011).

% --- BO quantum cosmology / derivation context ---
\bibitem{Brizuela2016}
D.~Brizuela, C.~Kiefer, and M.~Kr\"amer,
``Quantum-gravitational effects on gauge-invariant perturbations: de Sitter case,''
\emph{Phys.\ Rev.\ D} \textbf{93}, 104035 (2016).

\bibitem{Kamenshchik2018}
A.~Y.~Kamenshchik, A.~Tronconi, and G.~Venturi,
``The Born--Oppenheimer method, quantum gravity and inflation,''
\emph{Class.\ Quantum Grav.} \textbf{35}, 015012 (2018).

% --- Λ alternatives (sequestering / unimodular) ---
\bibitem{KaloperPadillaPRL2014}
N.~Kaloper and A.~Padilla,
``Sequestering the Standard Model Vacuum Energy,''
\emph{Phys.\ Rev.\ Lett.} \textbf{112}, 091304 (2014).

\bibitem{KaloperPadillaPRL2016}
N.~Kaloper, A.~Padilla, D.~Stefanyszyn, and G.~Zahariade,
``A Manifestly Local Theory of Vacuum Energy Sequestering,''
\emph{Phys.\ Rev.\ Lett.} \textbf{116}, 051302 (2016).

\bibitem{HenneauxTeitelboim1989}
M.~Henneaux and C.~Teitelboim,
``The Cosmological Constant and General Covariance,''
\emph{Phys.\ Lett.\ B} \textbf{222}, 195 (1989).

\bibitem{Ellis2011}
G.~F.~R.~Ellis, H.~van~Elst, J.~Murugan, and J.-P.~Uzan,
``On the Trace-Free Einstein Equations as a Viable Alternative to GR,''
\emph{Class.\ Quantum Grav.} \textbf{28}, 225007 (2011).

% --- DSI / log-oscillations context (optional module) ---
\bibitem{Sornette1998}
D.~Sornette,
``Discrete scale invariance and complex dimensions,''
\emph{Phys.\ Rep.} \textbf{297}, 239 (1998).

\bibitem{Zeng2019}
C.~Zeng, E.~D.~Kovetz, and M.~Kamionkowski,
``Searching for Oscillations in the Primordial Power Spectrum,''
\emph{Phys.\ Rev.\ D} \textbf{99}, 043517 (2019).

\end{thebibliography}

\end{document}